\section*{第 \ref{ch:b1:isostasy-gravity} 章——地壳均衡与重力异常}

\begin{solution}{exo:b1:isostasy-gravity:1}
(a) $\qty{-150}{\milli\gal} = -150\times\qty{e-5}{m.s^{-2}} = \qty{-1.5e-3}{m.s^{-2}}$，
且 $\num{1.5e-3}/9.81 = \num{1.53e-4}$（占 $g$ 的比例）。(b) $\qty{1}{\milli\gal} =
\qty{e-5}{m.s^{-2}}$，\num{1.02e-6}（占 $g$ 的比例）；$\qty{1}{\micro\gal} =
\qty{e-8}{m.s^{-2}}$，\num{1.02e-9}（占 $g$ 的比例）：能分辨一
微伽的重力仪可将重力测到约十亿分之一。
\end{solution}

\begin{solution}{exo:b1:isostasy-gravity:2}
(a) 赤道处两个正弦项均为零：$\gamma = \qty{9.780327}{m.s^{-2}}$。在
\ang{45}，$\sin^2\varphi = 0.5$ 且 $\sin^2 2\varphi = 1$：
$\gamma = 9.780327\times(1 + 0.0026512 - 0.0000058) = \qty{9.806200}{m.s^{-2}}$。
在两极，$\sin^2\varphi = 1$ 且 $\sin 2\varphi = 0$：
$\gamma = 9.780327\times1.0053024 = \qty{9.832186}{m.s^{-2}}$。
(b) $9.832186 - 9.780327 = \qty{0.051859}{m.s^{-2}}$，即
\qty{5186}{\milli\gal}。(c) $3392/5186 = 0.65$：离心效应解释了
约三分之二的差值；其余部分则由扁率解释，扁率使两极
更靠近地心。
\end{solution}

\begin{solution}{exo:b1:isostasy-gravity:3}
$r = h\rho_c/(\rho_m - \rho_c) = 3\times2800/500 = \qty{16.8}{km}$。\omterm{def:b1:layered-earth:layers}{地壳}厚度
增加了 $h + r = \qty{19.8}{km}$（即 $h\rho_m/(\rho_m - \rho_c) =
3\times3300/500$）。
\end{solution}

\begin{solution}{exo:b1:isostasy-gravity:4}
(a) $1/0.3086 = \qty{3.24}{m}$。(b) $0.01/0.3086 = \qty{0.032}{m}$：高度须
精确到约 \qty{3}{cm}，这就是测量时要
用卫星定位或水准测量记录每个测站高度的原因。
\end{solution}

\begin{solution}{exo:b1:isostasy-gravity:5}
(a) $\rho = 2700\times100/(100 + h)$：\qty{2673}{kg.m^{-3}}（对应 \qty{1}{km}），
\qty{2621}{kg.m^{-3}}（对应 \qty{3}{km}），\qty{2571}{kg.m^{-3}}（对应 \qty{5}{km}）。
(b) $\rho = (2700\times100 - 1030\times4)/96 = \num{265880}/96 =
\qty{2770}{kg.m^{-3}}$：海下的柱体密度更大。(c) 普拉特模型
要求在同一深度上存在横向密度变化。在洋底之下
确有这种变化：靠近洋脊的年轻、热的\omterm{def:b1:layered-earth:lithosphere}{岩石圈}比
远离洋脊的古老、冷的\omterm{def:b1:layered-earth:lithosphere}{岩石圈}更轻（见\cref{ch:b1:heat-budget}）。在大陆之下，
情况则相反，\omterm{def:b1:layered-earth:layers}{地壳}厚度随地形起伏而变化。
\end{solution}

\begin{solution}{exo:b1:isostasy-gravity:6}
(a) 测站 1：$g - \gamma = 9.807751 - 9.808911 = \qty{-0.001160}{m.s^{-2}} =
\qty{-116.0}{\milli\gal}$；\omterm{def:b1:isostasy-gravity:anomalies}{自由空气校正} $0.3086\times350 =
\qty{108.0}{\milli\gal}$，因此 $\Delta g_{\mathrm{FA}} = \qty{-8.0}{\milli\gal}$；
\omterm{def:b1:isostasy-gravity:anomalies}{布格校正} $0.112\times350 = \qty{39.2}{\milli\gal}$，因此
$\Delta g_{\mathrm B} = \qty{-47.2}{\milli\gal}$。
测站 2：$g - \gamma = \qty{-535.9}{\milli\gal}$；\omterm{def:b1:isostasy-gravity:anomalies}{自由空气校正}
$0.3086\times1850 = \qty{570.9}{\milli\gal}$，因此 $\Delta g_{\mathrm{FA}} =
\qty{+35.0}{\milli\gal}$；\omterm{def:b1:isostasy-gravity:anomalies}{布格校正} $0.112\times1850 =
\qty{207.2}{\milli\gal}$，因此 $\Delta g_{\mathrm B} = \qty{-172.2}{\milli\gal}$。
(b) 测站 2 处\omterm{def:b1:isostasy-gravity:anomalies}{布格异常}为强负值：山体之下
存在质量亏损，即轻质\omterm{def:b1:layered-earth:layers}{地壳}形成的山根。\omterm{def:b1:isostasy-gravity:anomalies}{自由空气异常}
与\omterm{def:b1:isostasy-gravity:anomalies}{布格校正}相比较小，表明亏损几乎与
山体相抵：二者在很大程度上已补偿（约为 $172/207$，
即 \qty{83}{\percent}，若板片模型对整个山脉成立）。
\end{solution}

\begin{solution}{exo:b1:isostasy-gravity:7}
(a) 与命题的证明相同，距离 $z$ 处的薄层
贡献为 $G\rho\,2\pi z\,h\int_0^a s\,(s^2 + z^2)^{-3/2}\,\mathrm ds$，此时
积分截止于 $s = a$：
$\int_0^a s\,(s^2 + z^2)^{-3/2}\,\mathrm ds = 1/z - 1/\sqrt{z^2 + a^2}$。由此
$g = 2\pi G\rho h\,\bigl(1 - z/\sqrt{z^2 + a^2}\bigr)$，当 $a \to \infty$ 时趋于板片
的值。(b) $z/\sqrt{z^2 + a^2} = 0.1$ 给出
$\sqrt{z^2 + a^2} = 10z$，$a = \sqrt{99}\,z = \qty{9.95}{km}$，对于
$z = \qty{1}{km}$。(c) 测站正下方约十倍埋深以内的岩石
提供了 \qty{90}{\percent} 的引力：地形校正
需要测站周围半径为几千米至几十千米范围内的地形
起伏，而不仅仅是其正下方的地形。
\end{solution}

\begin{solution}{exo:b1:isostasy-gravity:8}
(a) $r' = d(\rho_c - \rho_w)/(\rho_m - \rho_c) = 4\times1770/500 =
\qty{14.16}{km}$，\omterm{def:b1:layered-earth:layers}{地壳}厚度为 $35 - 4 - 14.16 = \qty{16.84}{km}$。
(b) 取单一\omterm{def:b1:layered-earth:layers}{地壳}密度的艾里模型预测洋壳厚度
比实际厚三倍以上。洋壳并非减薄的陆壳：它
是另一种更致密的岩石，地球的两个高度层级分属两类
\omterm{def:b1:layered-earth:layers}{地壳}，\omterm{def:b1:isostasy-gravity:hypsometry}{海陆起伏曲线}也表明了这一点。
\end{solution}

\begin{solution}{exo:b1:isostasy-gravity:9}
(a) 一角秒等于 $\pi/(180\times3600) = 1/\num{206265}$ 弧度：
\qty{5.236}{''} $= \qty{2.538e-5}{rad}$，而 \qty{15.885}{''} $=
\qty{7.701e-5}{rad}$。(b) $a_h = g\theta$：$9.79\times\num{2.538e-5} =
\qty{2.485e-4}{m.s^{-2}} = \qty{24.9}{\milli\gal}$（观测值），而
$9.79\times\num{7.701e-5} = \qty{75.4}{\milli\gal}$（由可见山体计算得出）。
(c) $1 - 5.236/15.885 = 0.670$：补偿抵消了可见山体引力的
约三分之二。
\end{solution}

\begin{solution}{exo:b1:isostasy-gravity:10}
(a) 在球外，球体对外的引力与位于球心的质点相同，
沿连线指向球心的引力为 $GM/(x^2 + z^2)$；竖直分量为其
乘以余弦 $z/\sqrt{x^2 + z^2}$，得 $GMz/(x^2 + z^2)^{3/2}$。
(b) $M = \tfrac43\pi(\qty{2000}{m})^3\times(\qty{-200}{kg.m^{-3}}) =
\qty{-6.70e12}{kg}$，球心正上方 $GM/z^2 = \num{6.674e-11}\times
(\num{-6.70e12})/(4000)^2 = \qty{-2.80e-5}{m.s^{-2}} = \qty{-2.80}{\milli\gal}$。
(c) $(1 + x^2/z^2)^{-3/2} = 1/2$ 可得 $x = z\sqrt{2^{2/3} - 1} = 0.766\,z =
\qty{3.07}{km}$。(d) 剖面在半极大处的半宽可给出
深度：$z = 1.305\,x_{1/2}$。埋深更大的异常体产生更宽、更平缓的异常。
\end{solution}

\begin{solution}{exo:b1:isostasy-gravity:11}
(a) $V = \tfrac13\pi R^2 H = \tfrac13\pi(\qty{15e3}{m})^2(\qty{3e3}{m}) =
\qty{7.07e11}{m^3}$，且 $\Delta m = (2800 - 1030)V = \qty{1.25e15}{kg}$。
(b) 质心位于海底以上 \qty{0.75}{km} 处，
即位于海面以下 $d = 5 - 0.75 = \qty{4.25}{km}$ 处。该处的位势异常
为 $T = -G\Delta m/d$，因此 $N = G\Delta m/(\gamma d) =
\num{6.674e-11}\times\num{1.25e15}/(9.81\times4250) = \qty{2.0}{m}$。
(c) 锥体的大部分质量位于其宽阔的下部，距
山顶正上方一点 5 至 \qty{16}{km}：整个质量分布上 $1/r$ 的平均值
小于 $1/\qty{4.25}{km}$，点质量高估了
位势。对薄圆盘求和，约为 \qty{1.1}{m}。
(d) 若其下存在轻质\omterm{def:b1:layered-earth:layers}{地壳}形成的山根，则相当于在更深处
设置同等质量的亏损；两者位势几乎相互抵消，隆起
要小得多，但并不为零，因为亏损比盈余离地表更远。
\end{solution}

\begin{solution}{exo:b1:isostasy-gravity:12}
(a) 隆升速率为 $-\mathrm dw/\mathrm dt = w_0\,e^{-t/\tau}/\tau = w/\tau$。
(b) $w = 818\,e^{-10000/4000} = 818\times0.0821 = \qty{67.1}{m}$，速率为
$\qty{67.1}{m}/\num{4000}\,\unit{\text{年}} = \qty{16.8}{mm.\text{年}^{-1}}$，超过实测值的
一倍半。(c) 当 $\tau = 3100$ 年时，$w =
818\,e^{-3.226} = \qty{32.5}{m}$，速率为 \qty{10.5}{mm.\text{年}^{-1}}；
当 $\tau = \num{68000}$ 年时，$w = 818\,e^{-0.147} = \qty{706}{m}$，速率
为 \qty{10.4}{mm.\text{年}^{-1}}。两种情形所隐含的剩余凹陷差异很大。
未补偿的凹陷即\omterm{def:b1:layered-earth:layers}{地幔}岩石的质量亏损：在大范围内，
它给出的\omterm{def:b1:isostasy-gravity:anomalies}{自由空气异常}约为 $-2\pi G\rho_m w$，即
$-0.0419\times3.3\times706 \approx \qty{-98}{\milli\gal}$（慢速情形），
而快速情形则为 \qty{-4}{\milli\gal}。重力测量（或
古海岸线的高度）可将二者区分开。(d) 时间常数
衡量\omterm{def:b1:layered-earth:layers}{地幔}在凹陷之下回流的快慢，即其动力黏度。
\end{solution}

\begin{solution}{pb:b1:isostasy-gravity:1}
\textbf{1.} $r = h\rho_c/(\rho_m - \rho_c) = 4\times2700/600 = \qty{18}{km}$。

\textbf{2.} \omterm{def:b1:layered-earth:layers}{地壳}厚 $35 + 4 + 18 = \qty{57}{km}$；其底面，即\omterm{def:b1:layered-earth:discontinuity}{莫霍面}，
位于海平面以下 $35 + 18 = \qty{53}{km}$。

\textbf{3.} 高原：\qty{57}{km} \omterm{def:b1:layered-earth:layers}{地壳}与 $60 - 53 = \qty{7}{km}$ \omterm{def:b1:layered-earth:layers}{地幔}，
$\num{57000}\times2700 + \num{7000}\times3300 = \num{1.539e8} + \num{2.31e7} =
\qty{1.770e8}{kg.m^{-2}}$。低地：$\num{35000}\times2700 + \num{25000}\times3300
= \num{9.45e7} + \num{8.25e7} = \qty{1.770e8}{kg.m^{-2}}$。两者相等，如
\cref{prop:b1:isostasy-gravity:equal-mass}所要求。

\textbf{4.} $p = g\times\text{单位面积质量} = 9.81\times\num{1.770e8} =
\qty{1.74e9}{Pa}$，约 \qty{1.7}{GPa}，两柱底部均为此值。

\textbf{5.} 山根须凭借较小的密度，补偿高原的质量。
每米山根移去 $\rho_m - \rho_c = \qty{600}{kg.m^{-3}}$ 的质量，而每米高原
增加 \qty{2700}{kg.m^{-3}} 的质量，故需 $2700/600 = 4.5$
米山根。若\omterm{def:b1:layered-earth:layers}{地壳}密度为 \qty{2900}{kg.m^{-3}}，则比值为
$2900/400 = 7.25$：对比度越小，山根越深。

\textbf{6.} \omterm{def:b1:layered-earth:discontinuity}{莫霍面}（即地震波速急剧增大的界面）将被发现
在高原边缘之下加深，从低地之下的 \qty{35}{km} 增至
高原之下的约 \qty{53}{km}。

\textbf{7.} $\rho = 2700\times100/104 = \qty{2596}{kg.m^{-3}}$。

\textbf{8.} 二者分别为 $\num{104000}\times2596.2 = \qty{2.700e8}{kg.m^{-2}}$
和 $\num{100000}\times2700 = \qty{2.700e8}{kg.m^{-2}}$。

\textbf{9.} $2700 - 2596 = \qty{104}{kg.m^{-3}}$，或 $104/2700 =
\qty{3.85}{\percent}$。

\textbf{10.} 由 $\Delta\rho/\rho = \alpha\,\Delta T$ 得 $\Delta T =
0.0385/\num{3e-5} = \qty{1280}{K}$，在整个 \qty{104}{km}
厚的柱体上的平均值。这将使柱体的大部分熔化：仅靠加热
不足以使该高原足够轻，普拉特机制在此不可能是主要机制。

\textbf{11.} 普拉特柱没有山根：\omterm{def:b1:layered-earth:discontinuity}{莫霍面}保持在约 \qty{35}{km}
高原之下，与低地之下相同。

\textbf{12.} \qty{52}{km} 接近艾里的 \qty{53}{km}，而远离普拉特的
\qty{35}{km}：该高原主要由\omterm{def:b1:layered-earth:layers}{地壳}山根补偿。

\textbf{13.} $0.0419\times2.7\times4000 = \qty{452.5}{\milli\gal}$。

\textbf{14.} 未补偿（见\cref{prop:b1:isostasy-gravity:broad}）：
$\Delta g_{\mathrm{FA}} \approx \qty{+452.5}{\milli\gal}$，
$\Delta g_{\mathrm B} \approx 0$。

\textbf{15.} 补偿时：$\Delta g_{\mathrm{FA}} \approx 0$，且
$\Delta g_{\mathrm B} \approx \qty{-452.5}{\milli\gal}$。山根所缺失的质量
给出 $-0.0419\times0.6\times\num{18000} = \qty{-452.5}{\milli\gal}$，结果相同，
因为 $\Delta\rho\,r = \rho_c h$。

\textbf{16.} $\sin^2\ang{30} = 0.25$ 且 $\sin^2\ang{60} = 0.75$：
$\gamma = 9.780327\times(1 + 0.0013256 - 0.0000044) = \qty{9.793249}{m.s^{-2}}$。
于是 $g - \gamma = 9.781085 - 9.793249 = \qty{-1216.4}{\milli\gal}$；\omterm{def:b1:isostasy-gravity:anomalies}{自由空气
校正}为 $0.3086\times4000 = \qty{1234.4}{\milli\gal}$，因此
$\Delta g_{\mathrm{FA}} = \qty{+18.0}{\milli\gal}$，且
$\Delta g_{\mathrm B} = 18.0 - 452.5 = \qty{-434.5}{\milli\gal}$。

\textbf{17.} $434.5/452.5 = 0.96$：高原约有 \qty{96}{\percent}
得到补偿（小的正\omterm{def:b1:isostasy-gravity:anomalies}{自由空气异常}即未补偿的
\qty{4}{\percent}）。

\textbf{18.} 不是。在宽广高原的中部，每一异常层都起板片的作用，
其引力不依赖于深度：只有单位面积的质量亏损起作用，
两种模型中它都相同，即 $\rho_c h$。两者的差别
仅出现在边缘附近，那里补偿质量的深度
改变了异常的形态。

\textbf{19.} 在两侧远离边缘处，\omterm{def:b1:isostasy-gravity:anomalies}{自由空气异常}接近于零。
在边缘以内数千米处，测点下方的高原距离很近，
且仅向一侧延伸很远：其引力仍接近其板片值的全部。
山根位于 \qty{35}{km} 或更深处，从测点看去是一个几乎笔直向下终止的物体：
其引力仅为其全部效应的一半多一点。两者
之差形成一个正值峰。在边缘以外数千米处，高原的吸引力
很小，而山根仍有接近一半的效应：形成一个负值
谷。剖面在每一边缘处呈正–负成对的形态，如
本章图所示。

\textbf{20.} \omterm{def:b1:layered-earth:layers}{地壳}现厚 $57 - 1 = \qty{56}{km}$，比低地
多 $\qty{21}{km}$。平衡时该超出量分为 $h' + r'$，其中
$r' = 4.5\,h'$：$h' = 21/5.5 = \qty{3.818}{km}$，$r' = \qty{17.18}{km}$。

\textbf{21.} 高原现高程为 \qty{3818}{m}。

\textbf{22.} 现位于地表的岩石原先位于旧地表以下 \qty{1}{km}，
即 \qty{3000}{m} 处；现位于 \qty{3818}{m}，已上升 \qty{818}{m}，
即 $1000\times\rho_c/\rho_m$。

\textbf{23.} 每侵蚀 1 千米，地表降低 $1 - 2700/3300 =
0.182$ 千米，因此要使地表降低 \qty{4}{km}，需
$4/0.182 = \qty{22}{km}$ 的侵蚀，即全部超出厚度 $h + r$。在古老而
被剥蚀的造山带中，地表岩石原先
深约二十千米，因此该造山带的深部山根已经出露。

\textbf{24.} 地表降低量为 $1000 - 818$：每侵蚀 1 千米岩石，
地表的净降低量为
\textbf{\qty[text-series-to-math=true]{182}{m}}。
\end{solution}
